\section{化学平衡・反応}

\subsection{令和7年度 問20}
\label{q:r07-reg-applied-20}

\noindent\textbf{原典：} sources/R07/17\_applied\_science.pdf、PDF物理7頁。

\textbf{公式正答：} 2（single）。

\par\textit{原典本文の校正と独自解説は未作成。分類のみの登録項目。}

\section{地球内部・プレート}

\subsection{令和6年度 問25}
\label{q:r06-reg-applied-25}

\noindent\textbf{原典：} sources/R06/17\_applied\_science.pdf、PDF物理8頁。

\textbf{公式正答：} 5（single）。

\par\textit{原典本文の校正と独自解説は未作成。分類のみの登録項目。}

\section{火山・マグマ}

\subsection{平成25年度 問25}
\label{q:h25-reg-applied-25}

\noindent\textbf{原典：} sources/H25/17\_applied\_science.pdf、PDF物理9頁。

\textbf{公式正答：} 3（single）。

\par\textit{原典本文の校正と独自解説は未作成。分類のみの登録項目。}

\subsection{平成25年度 問26}
\label{q:h25-reg-applied-26}

\noindent\textbf{原典：} sources/H25/17\_applied\_science.pdf、PDF物理10頁。

\textbf{公式正答：} 5（single）。

\par\textit{原典本文の校正と独自解説は未作成。分類のみの登録項目。}

\subsection{平成28年度 問28}
\label{q:h28-reg-applied-28}

\noindent\textbf{原典：} sources/H28/17\_applied\_science.pdf、PDF物理11頁。

\textbf{公式正答：} 5（single）。

\par\textit{原典本文の校正と独自解説は未作成。分類のみの登録項目。}

\subsection{平成28年度 問29}
\label{q:h28-reg-applied-29}

\noindent\textbf{原典：} sources/H28/17\_applied\_science.pdf、PDF物理12頁。

\textbf{公式正答：} 4（single）。

\par\textit{原典本文の校正と独自解説は未作成。分類のみの登録項目。}

\subsection{平成28年度 問35}
\label{q:h28-reg-applied-35}

\noindent\textbf{原典：} sources/H28/17\_applied\_science.pdf、PDF物理14頁。

\textbf{公式正答：} 2（single）。

\par\textit{原典本文の校正と独自解説は未作成。分類のみの登録項目。}

\subsection{平成29年度 問23}
\label{q:h29-reg-applied-23}

\noindent\textbf{原典：} sources/H29/17\_applied\_science.pdf、PDF物理8頁。

\textbf{公式正答：} 4（single）。

\par\textit{原典本文の校正と独自解説は未作成。分類のみの登録項目。}

\subsection{平成29年度 問29}
\label{q:h29-reg-applied-29}

\noindent\textbf{原典：} sources/H29/17\_applied\_science.pdf、PDF物理10頁。

\textbf{公式正答：} 3（single）。

\par\textit{原典本文の校正と独自解説は未作成。分類のみの登録項目。}

\subsection{平成30年度 問29}
\label{q:h30-reg-applied-29}

\noindent\textbf{原典：} sources/H30/17\_applied\_science.pdf、PDF物理10頁。

\textbf{公式正答：} 3（single）。

\par\textit{原典本文の校正と独自解説は未作成。分類のみの登録項目。}

\subsection{令和元年度 問23}
\label{q:r01-reg-applied-23}

\noindent\textbf{原典：} sources/R01/17\_applied\_science.pdf、PDF物理8頁。

\textbf{公式正答：} 2（single）。

\par\textit{原典本文の校正と独自解説は未作成。分類のみの登録項目。}

\subsection{令和元年度再試験 問22}
\label{q:r01-ret-applied-22}

\noindent\textbf{原典：} sources/R01\_retest/17\_applied\_science.pdf、PDF物理7頁。

\textbf{公式正答：} 4（single）。

\par\textit{原典本文の校正と独自解説は未作成。分類のみの登録項目。}

\subsection{令和元年度再試験 問23}
\label{q:r01-ret-applied-23}

\noindent\textbf{原典：} sources/R01\_retest/17\_applied\_science.pdf、PDF物理8頁。

\textbf{公式正答：} 3（single）。

\par\textit{原典本文の校正と独自解説は未作成。分類のみの登録項目。}

\subsection{令和元年度再試験 問24}
\label{q:r01-ret-applied-24}

\noindent\textbf{原典：} sources/R01\_retest/17\_applied\_science.pdf、PDF物理8頁。

\textbf{公式正答：} 5（single）。

\par\textit{原典本文の校正と独自解説は未作成。分類のみの登録項目。}

\subsection{令和3年度 問26}
\label{q:r03-reg-applied-26}

\noindent\textbf{原典：} sources/R03/17\_applied\_science.pdf、PDF物理9頁。

\textbf{公式正答：} 3（single）。

\par\textit{原典本文の校正と独自解説は未作成。分類のみの登録項目。}

\subsection{令和5年度 問25}
\label{q:r05-reg-applied-25}

\noindent\textbf{原典：} sources/R05/17\_applied\_science.pdf、PDF物理9頁。

\textbf{公式正答：} 3（single）。

\par\textit{原典本文の校正と独自解説は未作成。分類のみの登録項目。}

\subsection{令和7年度 問26}
\label{q:r07-reg-applied-26}

\noindent\textbf{原典：} sources/R07/17\_applied\_science.pdf、PDF物理10頁。

\textbf{公式正答：} 2（single）。

\par\textit{原典本文の校正と独自解説は未作成。分類のみの登録項目。}

\section{火砕密度流の形成と地形に沿う分布}

\subsection{令和元年度 問25}
\label{q:r01-reg-applied-25}

\noindent\textbf{原典：} sources/R01/17\_applied\_science.pdf、PDF物理9頁。

\textbf{公式正答：} 5（single）。

\par\textbf{検討未確定：} 大規模火砕流/プリニー/カルデラの一般化とプレー型の用語範囲は全類型の一次資料照合が未完了である。公式/独立5と確認済みの非単純分布を保持し部分校正に残す。

\paragraph{問題の要旨} 火砕流の発生・混合物・堆積性状と、地形を反映する分布を比較する。

\Needspace{4\baselineskip}\paragraph{確認した条件・記述（独自要約）}\mbox{}\par

\noindent\textbf{判断方向：} 最も不適切・誤りを選ぶ\par

\noindent\textbf{原典照合：} 15原典画像で18問の問題幹、全5肢、否定、専門語、化学式、候補順を確認。\par

\noindent\textbf{適用条件：} 流れと構造の典型対応は一意成因を保証しない。用語、時間間隙、保存、試料処理の条件を個別に示す。\par

\paragraph{問われる核心} 火砕流の発生・混合物・堆積性状と、地形を反映する分布を比較する。

\paragraph{必要知識} 火砕密度流の粒子と気体、ドーム崩落、噴煙柱、地形と分布、降下と流下の区別。

\paragraph{考え方・正答を選ぶ理由（査読待ち初稿）} 火砕流は地面に沿う密度流で、経路と堆積域は谷や尾根、流れの濃度に依存する。USGSには重なった扇状・舌状・ローブ状の分布及びメラピの谷埋積と希薄サージ分離の報告がある。従って単純な長円形が一般的だとする5が最も不適切である。降下火山灰は風に運ばれるため別の分布モデルを使う。大規模噴出・カルデラやプレー型という呼称は、この資料確認だけで全例を普遍化せず原典の典型説明として保持する。

\Needspace{4\baselineskip}\paragraph{選択肢の検討（査読待ち初稿）}\mbox{}\par

\noindent\textbf{1：適切（条件を含めて検討）} 一般に分級が悪く塊状となる。多粒径を混合して急速に堆積する流れの代表的特徴である。希薄サージ等の全堆積物へ塊状を強制しない。\par

\noindent\textbf{2：適切（条件を含めて検討）} 大規模火砕流は高発泡のプリニー式に伴い通常カルデラを作る。高発泡マグマの大規模噴出とカルデラ形成の典型を述べるが、全活動様式の必要十分条件としては確認していない。\par

\noindent\textbf{3：適切（条件を含めて検討）} 高温火砕物・火山ガス・取り込む大気が地面に沿って流れる。高温粒子と気体の混合物が地形に沿って流下する現象に対応する。固体だけ又は溶岩だけではない。\par

\noindent\textbf{4：適切（条件を含めて検討）} ドームの爆発型をプレー型、重力崩落型をメラピ型とする。メラピのドーム崩落起源は一次論文で確認した。プレー型という原典呼称と全類型の整理は追加確認を要する。\par

\noindent\textbf{5：誤り・最も不適切} 分布は火口を長軸の一端に持つ単純な長円が多い。谷沿いの舌状・ローブ状、複数谷の埋積などを取り得る。風で伸びる降下物の模式分布と混同する一般化である。\par

\paragraph{再利用する考え方} 観察できる構造と成因・環境の解釈を区別する。代表モデルに適用条件を添え、単一の形や測定値から唯一の環境を断定しない。

\paragraph{共通理論} \hyperref[theory:volcanology:pyroflow]{第\ref*{theory:volcanology:pyroflow}節}, \hyperref[theory:volcanology:source]{第\ref*{theory:volcanology:source}節}, \hyperref[theory:volcanology:distribution]{第\ref*{theory:volcanology:distribution}節}

\par 参考：\cite{volc_usgs_merapi,volc_usgs_msh_flow}。

\par 難易度 3（校準初稿）、優先度未評価。

\section{未確定}

\subsection{令和2年度 問24}
\label{q:r02-reg-applied-24}

\noindent\textbf{原典：} sources/R02/17\_applied\_science.pdf、PDF物理9頁。

\textbf{公式正答：} 3（single）。

\par\textit{原典本文の校正と独自解説は未作成。分類のみの登録項目。}
